\section{Primary Contribution --- The Crossover Point}

The primary contribution of this thesis is the identification and measurement of a
\textbf{crossover point}, formally defined as:

\begin{quote}
The concurrent workload level at which \textbf{(1)} the \vm{} deployment's job success rate
drops below 95\% due to resource contention, \textbf{AND (2)} the Kubernetes deployment's
total execution time including pod scheduling overhead and container startup latency becomes
lower than the \vm{}'s execution time under active contention.
\end{quote}

This crossover point is not an interpretation --- it is a computable result derived directly
from experimental data. It gives engineers a specific, data-backed threshold for when the
Kubernetes migration becomes net beneficial.
